\begingroup
\setlength{\emergencystretch}{2em}
\subsection*{Архив: архитектура и организация компьютера (x86-64)}
\label{lecture04-journal}
Сквозной пример --- разработка журнала оценок группы. Нужно вести оценки
16 студентов по четырём работам на компьютере, где \emph{нет гостевой
операционной системы}. Презентация и исполняемые контрольные версии
находятся в \path{slides/lecture04-journal/} и
\path{demos/lecture04-journal/}. Семинар №4 закрепляет адреса и память отдельно.

Учебная машина --- QEMU \texttt{pc} (i440FX) с 64 MiB RAM и CPU x86-64.
У неё VGA-текстовый экран, PS/2-клавиатура, совместимый IDE-контроллер
и отдельный диск 1 MiB.
BIOS и GRUB \emph{присутствуют}: прошивка передаёт управление загрузчику,
тот загружает нашу программу. После передачи управления нет Linux и нельзя
автоматически использовать \texttt{printf}, \texttt{malloc}, системные вызовы
и файлы с именами. Компилятор и GRUB работают \emph{при подготовке образа}
в Docker, а не внутри загруженного учебного компьютера.

\begin{center}
\small
\begin{tabular}{|p{0.21\linewidth}|p{0.62\linewidth}|}
\hline
Минуты & Потребность и архитектурный вопрос \\
\hline
0--6 & Обзор ПК, ограничения машины без гостевой ОС \\
6--18 & Запуск и один символ: CPU, RAM, регистры, VGA \\
18--30 & Таблица студентов: байты, адреса, массивы, размеры \\
30--44 & Ввод: PS/2, шины, опрос, IRQ1 и очередь \\
44--58 & Средние: АЛУ, код и данные через кэши, локальность \\
58--64 & Стек, статический резерв и понятие кучи \\
64--80 & Сохранение: RAM, HDD/SSD, ATA PIO; сравнение с DMA \\
80--87 & GPU и современные соединения; зачем нужна ОС \\
87--90 & Итоги и самопроверка \\
\hline
\end{tabular}
\end{center}

\paragraph{От BIOS до нашей точки входа.}
В лабораторной используется QEMU PC в BIOS-совместимом режиме. После
сброса BIOS выполняет POST и базовую инициализацию, затем выбирает
загрузочный CD-ROM. По El Torito BIOS читает стартовый образ GRUB с ISO
в память и передаёт ему управление. Это ещё не загрузка ядра журнала:
прошивка не разбирает его C-код и не вызывает \texttt{kmain}.

GRUB читает конфигурацию, где указано ядро Multiboot2, загружает
\texttt{kernel.elf}, размещает его ELF-сегменты в RAM и формирует сведения
Multiboot2. Затем он передаёт управление 32-битной точке входа \texttt{\_start}.
Фрагмент конфигурации:
\begin{lstlisting}[basicstyle=\ttfamily\small]
menuentry "Course journal (bare metal)" {
    multiboot2 /boot/kernel.elf
    boot
}
\end{lstlisting}
Это конкретный путь BIOS/El Torito в QEMU PC (i440FX), а не UEFI.

\paragraph{Стартовый код: подготовить процессор для C.}
При входе от GRUB процессор находится в 32-битном protected mode.
Стартовый код создаёт PML4, PDPT и таблицу страниц с тождественным
отображением первых 1 GiB крупными страницами по 2 MiB. В одном каталоге
512 записей, каждая по 8 байтов. Затем включает PAE, задаёт CR3, выставляет
EFER.LME, загружает GDT, включает paging и делает дальний переход
в 64-битный сегмент. Основные регистровые шаги из \texttt{boot.asm}:
\begin{lstlisting}[basicstyle=\ttfamily\small]
mov eax, cr4
or eax, 1 << 5          ; PAE
mov cr4, eax
mov eax, pml4
mov cr3, eax
mov ecx, 0xc0000080     ; EFER
rdmsr
or eax, 1 << 8          ; LME
wrmsr
\end{lstlisting}
Перед вызовом C задаются сегментные регистры и стек:
\begin{lstlisting}[basicstyle=\ttfamily\small]
long_start:
    mov ax, 0x10
    mov ds, ax
    mov es, ax
    mov ss, ax
    mov rsp, stack_top
    xor rbp, rbp
    call kmain
\end{lstlisting}
BIOS здесь уже завершил свою часть; GRUB передал ядру стартовые сведения,
а подготовку long mode выполняет NASM. Детали заголовка Multiboot2,
таблиц страниц и GDT можно разобрать дополнительно.

\paragraph{Версия 1: получить первый результат.}
Нам нужна одна видимая оценка. Но кто начнёт выполнять команды?
GRUB передал Multiboot2-загрузку стартовой точке \texttt{\_start};
NASM подготовил процессор и стек, после чего вызвал \texttt{kmain} на C.
Теперь добавим самую простую полезную функцию --- вывод одной записи.

В первой версии C-код пишет символы в текстовый буфер VGA по адресу
\texttt{0xB8000}. Каждая позиция --- два байта: код символа и атрибут цвета.
Для краткости на экране \texttt{S01 grade 07}. Эти записи относятся
к памяти устройства, а не к обычному массиву DRAM. Отдельный последовательный
порт COM1 отправляет \texttt{BOOT STAGE=1} и \texttt{READY}: так тест
проверяет реальную загрузку без чтения пикселей экрана. Кодирование VGA
не является поддержкой Unicode; поэтому интерфейс учебной машины использует
латиницу, а объяснения в пособии и презентации --- русский язык.
Фрагмент реального \texttt{kernel.c} сначала ограничивает позицию экраном,
затем помещает цвет и код символа в одно 16-битное слово:
\begin{lstlisting}[language=C,basicstyle=\ttfamily\small]
static void vtext(unsigned row, unsigned col,
                  const char *s, uint8_t color) {
    while (*s && row < 25 && col < 80)
        VGA[row * 80 + col++] =
            ((uint16_t)color << 8) | (uint8_t)*s++;
}
/* redraw(), stage 1 */
vtext(4, 2, "S01  grade 07", 0x0f);
\end{lstlisting}
Это конкретное аппаратное отображение, а не переносимая функция вывода.
Версия 1 вызывает её с координатами и строкой
\texttt{"S01  grade 07"}. Индекс \texttt{row * 80 + col} указывает
позицию 80\,$\times$\,25 текстового экрана; одно присваивание записывает
код символа и цвет в 16-битное слово.

\paragraph{Версия 2: места для всей группы.}
Теперь у каждого из 16 условных студентов \texttt{S01}--\texttt{S16}
по четыре работы \texttt{W1}--\texttt{W4}. Массив
\texttt{uint8\_t grade[16][4]} занимает ровно $16\cdot4=64$ байта,
потому что каждый элемент имеет размер один байт и элементы идут подряд.
Разрешённые оценки --- $0,1,\ldots,10$; \texttt{0xFF} означает
«не выставлена». Поэтому \texttt{0} нельзя использовать как признак
отсутствия оценки. Для допустимых индексов байтовое смещение равно
$4i+j$, где $i$ --- номер студента с нуля, $j$ --- номер работы.
Так, \texttt{grade[1][2]} лежит на смещении $4\cdot1+2=6$ от начала.
Стадии 2--6 объявляют эту таблицу и sentinel:
\begin{lstlisting}[language=C,basicstyle=\ttfamily\small]
#define STUDENTS 16
#define WORKS 4
#define UNSET 0xffu
static uint8_t grade[STUDENTS][WORKS];
\end{lstlisting}
При старте \texttt{kmain} выполняет:
\begin{lstlisting}[language=C,basicstyle=\ttfamily\small]
for (unsigned i = 0; i < STUDENTS; i++)
    for (unsigned j = 0; j < WORKS; j++)
        grade[i][j] = UNSET;
\end{lstlisting}
По правилам C статические объекты изначально имеют нулевое значение;
загрузчик ELF размещает BSS. Явный цикл нужен не для «исправления случайной
RAM», а чтобы записать в каждую ячейку выбранное нами значение «не выставлено»:
ноль уже зарезервирован под оценку.
Отрисовка выбирает поле теми же индексами:
\begin{lstlisting}[language=C,basicstyle=\ttfamily\small]
uint8_t g = grade[i][j];
if (g == UNSET) vtext(row, col, "--", color);
else number2(row, col, g, color);
\end{lstlisting}
Цвет выделения и позиции \texttt{row}, \texttt{col} рассчитываются
в окружающем цикле \texttt{redraw}.

Адрес буфера --- не его содержимое. Наша C-программа использует
виртуальные адреса, пока стартовый код отображает их тождественно для
первых 1 GiB. Это выбранная конфигурация, а не свойство языка C.
Условный рисунок с базой \texttt{0x1000} и реальные адреса не смешиваем.
Массив и вывод в VGA тоже различны: отрисовка лишь отображает копию
данных на экране, не делает дисплей долговременным носителем оценок.

\paragraph{Версия 3: ввод посредством опроса.}
Клавиатура PS/2 отправляет скан-коды контроллеру i8042. CPU проверяет
младший бит порта состояния \texttt{0x64}: когда данные готовы,
читает скан-код из \texttt{0x60}. Это портовый ввод-вывод \texttt{in},
а не обычное чтение C-указателем из RAM. Программа последовательно
проверяет статус (polling); пока клавиш нет, ядро занято ожиданием.
Версия 3 выполняет именно такой опрос в главном цикле:
\begin{lstlisting}[language=C,basicstyle=\ttfamily\small]
if (in8(0x64) & 1)
    handle(in8(0x60));
\end{lstlisting}
Сначала читается status register контроллера, затем data register.

Стрелки выбирают ячейку; клавиши \texttt{0}--\texttt{9} присваивают оценки,
\texttt{A} задаёт 10, Backspace очищает. Изменение выполняется сразу
по нажатию; \texttt{Enter} и \texttt{Escape} в этой учебной версии
не используются. Скан-код --- не готовый символ: таблица конкретных
make-кодов клавиш переводит его в действие интерфейса. Break-коды
отпускания не означают ещё одного ввода оценки. Реальные USB-клавиатуры
используют другой контроллер и протокол; это уже тема сравнения в конце.
В обработчике make-кодов код выбранной оценки попадает в рабочую таблицу:
\begin{lstlisting}[language=C,basicstyle=\ttfamily\small]
if (scan == 0x1e) value = 10;      /* A */
else if (scan == 0x0e) value = UNSET; /* Backspace */
else return;
grade[selected_student][selected_work] = (uint8_t)value;
report_grade();
redraw();
\end{lstlisting}
Полная функция \texttt{handle} сначала разбирает отпускание клавиш,
префикс стрелок \texttt{E0} и цифры 0--9.

\paragraph{Версия 4: вместо постоянного опроса --- IRQ1.}
Контроллер прерываний PIC перенаправляет аппаратный IRQ1 на вектор 33.
Стартовый код устанавливает запись IDT, разрешает только клавиатурный
IRQ, а таймер IRQ0 оставляет замаскированным. Обработчик на NASM
сохраняет регистры, вызывает короткую C-функцию, которая считывает
скан-код и кладёт его в кольцевую очередь, посылает подтверждение PIC
и возвращается через \texttt{iretq}. Основной цикл выбирает события
из очереди, проверяет их и меняет таблицу; отрисовка не выполняется
внутри обработчика. При переполнении 64-элементной очереди новое событие
теряется: это известное ограничение учебного интерфейса.
Функция \texttt{keyboard\_irq} только читает готовый байт,
помещает его в очередь и посылает PIC подтверждение EOI:
\begin{lstlisting}[language=C,basicstyle=\ttfamily\small]
void keyboard_irq(void) {
    if (in8(0x64) & 1)
        enqueue(in8(0x60));
    out8(0x20, 0x20);       /* PIC: EOI */
}
\end{lstlisting}
Запись IDT создаёт gate на адрес обработчика и загружает IDTR:
\begin{lstlisting}[language=C,basicstyle=\ttfamily\small]
uintptr_t p = (uintptr_t)keyboard_entry;
idt[33] = (struct idt_gate){
    (uint16_t)p, 8, 0, 0x8e,
    (uint16_t)(p >> 16), (uint32_t)(p >> 32), 0
};
struct idtr desc = { sizeof idt - 1, (uintptr_t)idt };
__asm__ volatile ("lidt %0" : : "m"(desc));
\end{lstlisting}
Затем PIC перенаправляется и разрешается лишь IRQ1; IRQ0 таймера остаётся
замаскированным. Буфер PS/2 очищается до \texttt{sti}. NASM-обёртка
сохраняет регистры, выравнивает стек для вызова C, восстанавливает состояние
и завершает обработку через \texttt{iretq}. Основной цикл версии 4 выбирает
готовые скан-коды из очереди и лишь там вызывает \texttt{handle} и отрисовку.

Прерывание --- уведомление о событии и передача управления, а не
сама оценка в памяти и не обязательное переключение процесса.
Отдельная таблица в IDT нужна потому, что в long mode формат дескриптора
иной, чем в 32-битном режиме. Сравните эту процедуру с логической схемой
обработки прерываний основной лекции (с.~\pageref{lecture04-interrupts}).

\paragraph{Версия 5: посчитать среднее.}
Перебираем только выставленные оценки; для ряда \texttt{0, 10, --, --}
сумма 10, количество 2 и среднее 5.00. Если оценок нет, выводим
\texttt{--.--}, а не делим на ноль. Для двух десятичных знаков используем
целочисленные сотые с округлением половины вверх:
\[
\operatorname{hundredths}
=\left\lfloor\frac{100\cdot\operatorname{sum}+\lfloor n/2\rfloor}{n}\right\rfloor,
\qquad n>0.
\]
Например, оценки 0, 2, 10 дают $1200/3=400$ сотых, то есть 4.00;
оценки 0, 1, 1 дают $(200+1)/3=67$ сотых, то есть 0.67.
Диапазон мал, переполнения используемых 32-битных сумм нет.
Реальный код проходит по четырём работам и исключает только sentinel:
\begin{lstlisting}[language=C,basicstyle=\ttfamily\small]
unsigned sum = 0, count = 0;
for (unsigned j = 0; j < WORKS; j++) {
    uint8_t g = grade[i][j];
    if (g != UNSET) { sum += g; count++; }
}
if (count)
    hundredths = (sum * 100 + count / 2) / count;
\end{lstlisting}
Отображение суммы тоже производится кодом: целая часть и две цифры сотых
записываются через \texttt{number2} и \texttt{vchar} в VGA.

Теперь выгодно проследить два потока. Байты инструкций процедуры берутся
из памяти через подсистему выборки (часто L1I), байты последовательных
оценок --- через подсистему данных (часто L1D); АЛУ выполняет сложение,
регистры удерживают промежуточные результаты. Условная линия кэша
64 байта могла бы вместить все 64 оценки, если начало массива выровнено,
но не обязана: массив может пересекать границу линий. L1-промах не означает
обязательное обращение к DRAM. Мы не измеряем настоящие задержки кэша
в QEMU: на рисунке показываем объяснительную модель.

\paragraph{Стек и размер журнала.}
Возвраты из функций, часть локальных значений и сохранённые регистры
требуют рабочего пространства. Стек объёмом 16 KiB подготовлен заранее;
локальные C-переменные не обязаны все буквально находиться в нём.
Таблица выделена статически: число записей известно заранее.
Если бы список студентов рос во время исполнения, понадобился бы
аллокатор для динамических блоков и политика освобождения. Без ОС
\texttt{malloc} не появится сам по себе; его можно реализовать поверх
имеющейся памяти, но в этом занятии он не требуется. Память прошивки,
RAM, кэши, регистры и накопитель не являются взаимозаменяемыми ресурсами.

\paragraph{Версия 6: оценки должны пережить выключение.}
RAM энергозависима. Клавиша \texttt{S} сериализует таблицу в один сектор
512 байтов отдельного диска и записывает через ATA PIO по адресу блока
LBA 1. Чтение при старте и команда \texttt{L} выполняют обратную операцию.
В этом учебном протоколе CPU пересылает все 256 слов сектора отдельными
инструкциями \texttt{inw}/\texttt{outw} через порт \texttt{0x1F0}.
Отдельный flush дискового кэша запрашивается до сообщения \texttt{SAVED}.
Это не вызов файловой системы: имён файлов и каталогов у учебного журнала
нет. При запуске проверяем сигнатуру, версию, размеры, значения оценок и
контрольную сумму до изменения таблицы в RAM.
В \texttt{save} сектор сначала заполняется нулями, затем поля кодируются
по определённым смещениям. Фрагмент цикла оценок и записи checksum:
\begin{lstlisting}[language=C,basicstyle=\ttfamily\small]
for (unsigned i = 0; i < 512; i++) sector[i] = 0;
sector[0] = 'G'; sector[1] = 'R';
sector[2] = 'D'; sector[3] = '4';
sector[4] = 1; sector[5] = STUDENTS;
sector[6] = WORKS; sector[7] = 0;
for (unsigned i = 0; i < STUDENTS; i++)
    for (unsigned j = 0; j < WORKS; j++)
        sector[8 + i * WORKS + j] = grade[i][j];
uint32_t c = checksum();
for (unsigned i = 0; i < 4; i++)
    sector[72 + i] = (uint8_t)(c >> (8 * i));
\end{lstlisting}
Далее \texttt{disk\_io(1)} выдаёт команду ATA, ждёт готовность и циклом
\texttt{outw} переносит 256 слов через порт. Это PIO: CPU сам переносит
каждое слово; DMA был бы другим вариантом работы контроллера.
Цикл передачи байтового буфера на порт выглядит так:
\begin{lstlisting}[language=C,basicstyle=\ttfamily\small]
for (unsigned i = 0; i < 256; i++) {
    if (write)
        out16(0x1f0, (uint16_t)(sector[i*2] |
            ((uint16_t)sector[i*2+1] << 8)));
    else {
        uint16_t w = in16(0x1f0);
        sector[i*2] = (uint8_t)w;
        sector[i*2+1] = (uint8_t)(w >> 8);
    }
}
\end{lstlisting}
Ветка чтения использует тот же порт, но принимает слова от контроллера.

\begin{center}
\small
\begin{tabular}{|l|p{0.6\linewidth}|}
\hline
Смещение & Байты сектора \\
\hline
0--3 & \texttt{GRD4} (ASCII) \\
4--7 & версия 1; числа 16 и 4; зарезервированный ноль \\
8--71 & оценки 0--10 либо \texttt{FF} в порядке студентов и работ \\
72--75 & 32-битная FNV-1a от первых 72 байтов, little-endian \\
76--511 & нули \\
\hline
\end{tabular}
\end{center}

Формат не получается автоматически из размещения структуры C: там
могут быть padding и иные размеры полей. Контрольная сумма выявляет
случайное повреждение, но не защищает от намеренной подмены. Ошибки
\texttt{EMPTY}, \texttt{NO\_DISK}, \texttt{BAD\_FORMAT},
\texttt{UNSUPPORTED\_VERSION}, \texttt{BAD\_CHECKSUM} и
\texttt{BAD\_GRADE} сохраняют текущую таблицу без изменений.
В \texttt{load} проверка оценок предшествует отдельному циклу копирования:
\begin{lstlisting}[language=C,basicstyle=\ttfamily\small]
for (unsigned i = 0; i < STUDENTS * WORKS; i++)
    if (sector[8+i] != UNSET && sector[8+i] > 10) {
        set_notice("BAD_GRADE"); return;
    }
for (unsigned i = 0; i < STUDENTS; i++)
    for (unsigned j = 0; j < WORKS; j++)
        grade[i][j] = sector[8 + i * WORKS + j];
\end{lstlisting}
Одна перезаписываемая копия сектора не гарантирует целостность при
внезапном отключении питания прямо в процессе записи; реальной системе
нужна более надёжная схема обновления.

\paragraph{Где в этом примере DMA и GPU?}
На учебном IDE-контроллере использовано PIO. При \emph{DMA} ОС или
низкоуровневый код настроили бы буфер, после чего контроллер перенёс
бы блок в RAM или из RAM без побайтовой пересылки CPU. Прерывание могло
бы сообщить о завершении: оно не тождественно переносу данных. Но это
сравнение, не скрытая функция нашего кода. Так же наш VGA-текстовый
буфер не выполняет параллельные вычисления на GPU: современный дискретный
GPU, VRAM и пути через PCIe --- другие механизмы, требующие иных драйверов.
HDD ждёт движение головки и вращение; SSD использует NAND Flash и
контроллер. Даже если оба подключены к системе через контроллеры, их
задержки и пропускная способность различаются.

Современные USB, PCIe, SATA и NVMe отвечают за разные соединения и
протоколы; M.2 задаёт конструктивное исполнение и разъём накопителя.
Шина связывает устройства и контроллеры обменом адресов, данных и
управляющей информации, но у современного ПК нет одной общей шины
для всего перечисленного. На реальной машине для регулярного ввода,
вывода, сохранения, работы с процессами и правами удобнее ОС. Её
назначение и запуск обычной программы изучаются на лекции №5.

\paragraph{Самопроверка.}
Почему BIOS/GRUB не являются гостевой ОС нашего журнала? Чем отсутствие
оценки отличается от нуля? Почему \texttt{grade[1][2]} имеет смещение 6?
Почему очередь нужна между обработчиком IRQ и основной программой?
Когда обращение к кэшу инструкций отличается от чтения оценки?
Что при записи сектора делает CPU в PIO, а что делал бы контроллер в DMA?
Почему \texttt{SAVED} не означает появление файла в каталоге?

\textbf{Ответы.} После передачи управления они не предоставляют
журналу сервисов ОС; \texttt{FF} выделен под пропуск; четыре работы на
студента, поэтому $1\cdot4+2=6$; обработчик быстро фиксирует событие,
а основная программа меняет таблицу; команды и операнды --- разные потоки;
при PIO CPU переносит слова через порт, при DMA контроллер передаёт
блок; мы записали адресуемый сектор без файловой системы.

\paragraph{Как воспроизвести.}
Из корня репозитория выполните нужную команду:
\begin{verbatim}
make lecture4-journal-check
make lecture4-journal-run
\end{verbatim}
Первая собирает шесть версий и проверяет их в QEMU. Вторая запускает
шестую версию; noVNC доступен по адресу
\url{http://localhost:6080/vnc.html}.
Диск \path{build/lecture04-journal/grades.img} сохраняется между
запусками. Клавиши и формат описаны в
\path{demos/lecture04-journal/README.md}.
\endgroup
